\subsection{An actual quotient lattice and the prescribed CM lifting fibre}

\begin{result}{E.2.13}{A fixed-\(\Q\) elliptic quotient outside the admissible closure}{Proved}
Let
\[
 E:\ y^2=x(x-1)(x-3),\qquad
 B:\ Y^2=X^3+8X^2+4X .
\]
Let \(\mathcal C_{\mathrm{adm}}(E)\) be the full rigid exact tensor
category generated by the actual integral Hodge structure
\(H_1(E(\C),\Z)\), using finite sums, tensor products, duals,
saturated Hodge subobjects and free Hodge quotients. Its morphisms
are all integral Hodge morphisms between the objects so obtained;
direct sums of pure weights carry their weight grading.
The actual \(H_1(B(\C),\Z)\) has no admissible presentation in this
category. An explicit integral tensor automorphism on
\(\mathcal C_{\mathrm{adm}}(E)\) fails to extend naturally across
the \(\Q\)-defined isogeny \(q:E\to B\).
\end{result}
\begin{proof}
The discriminants of the displayed equations are \(576\) and
\(12288\), respectively. Define
\[
 q(x,y)=
 \left(x-4+\frac3x,\ y\left(1-\frac3{x^2}\right)\right).
 \tag{E.2.13a}
\]
Substitution verifies the equation of \(B\). The rational map extends
between the smooth projective curves, takes the origin to the origin,
and satisfies
\[
 q^*(dX/Y)=dx/y.
\]
In the corresponding normalized complex universal covers its lift
therefore has derivative one and takes zero to zero, so it is the
map \(z\mapsto z\). This proves directly that \(q\) is a group
homomorphism. The nontrivial involution
\[
 (x,y)\longmapsto(3/x,-3y/x^2)
\]
fixes both coordinates in (E.2.13a); it is translation by
\(P=(0,0)\). Conversely \(x\) satisfies
\(x^2-(X+4)x+3=0\), and the displayed formula for \(Y\) then
recovers \(y\) generically. Thus \(q\) has degree exactly two and
kernel \(\{O,P\}\).

The curve \(E\) has
\[
 j(E)=\frac{21952}{9}.
\]
This rational number is not an algebraic integer; CM
\(j\)-invariants are algebraic integers by
\cite[Theorem 12.14 and Corollaries 12.15--12.16,
pp.119--120]{EllMilneMF}. Hence \(E\) has no CM. Its
Mumford--Tate group on \(H_1(E,\Q)\) is \(\GL_2\), by the group
calculation in A.4.9 and
\cite[Proposition 5.14]{GenKSV2026}.

Every \(T\in\mathrm{SL}(H_1(E,\Z))\) consequently defines an
integral tensor automorphism of the forgetful functor on
\(\mathcal C_{\mathrm{adm}}(E)\). To verify every lattice operation,
first act by the usual tensor representation on each integral
tensor word. Every rational Hodge subspace is
Mumford--Tate stable, and every rational Hodge morphism commutes
with this action. A saturated subobject \(N\subset P\) has
\[
 N=N_{\Q}\cap P.
\]
Therefore both \(T\) and \(T^{-1}\) preserve \(N\), and induce
mutually inverse maps on every free quotient. Induction constructs
the maps on all admissible expressions. The rational
Mumford--Tate action proves naturality for all original integral
Hodge morphisms and for isomorphisms between expressions, so the
construction is independent of the presentation. Tensor, unit and
dual compatibilities follow from the same representation.
The determinant-one condition fixes the labelled Tate object
and the elliptic polarization.

Choose a lattice basis \(e_1,e_2\) of
\(\Lambda_E=H_1(E(\C),\Z)\) in which \(P\) is represented by
\(e_1/2\). A nonzero vector modulo two has a primitive integral
lift, and such a lift extends to a basis, so this choice is possible.
The analytic quotient and the normalization \(dq=1\) identify
\[
 \Lambda_B=\Lambda_E+\Z(e_1/2).
\]
In the basis \(\mu_1=e_1/2,\ \mu_2=e_2\) of \(\Lambda_B\), the
matrix of \(q_*\) is
\[
 Q=\begin{pmatrix}2&0\\0&1\end{pmatrix}.
\]
Use the tensor automorphism constructed above with
\[
 T_E=J=\begin{pmatrix}0&-1\\1&0\end{pmatrix}.
\]
Naturality for \(q_*\) would require \(T_BQ=QJ\), and therefore
\[
 T_B=QJQ^{-1}
    =\begin{pmatrix}0&-2\\1/2&0\end{pmatrix}.
 \tag{E.2.13b}
\]
This matrix is not integral. Thus the stated extension does not
exist. An admissible presentation of the actual \(H_1(B,\Z)\)
would produce such an extension: fullness includes the integral
Hodge morphism \(q_*\), to which the constructed tensor
automorphism must be natural. This contradiction proves the
nonmembership assertion.

The same obstruction applies to a motivic admissible presentation
with these specified actual Betti lattices. Its Betti--Hodge
realization is exactly an admissible presentation of the sort just
excluded, and the isogeny supplies the algebraic correspondence
\(q_*\). Passing to cohomology uses the inverse dual
\(T=(\eta^{-1})^\vee\), so it does not change the obstruction.
The identities of curves, differentials, discriminants and matrices
are checked exactly by \texttt{verify\_numberfield\_descent.py}.

This supplies an actual failure of the additional presentation
step needed to identify E.2.12 with the full geometric diagram,
over the one fixed field \(\Q\). Enlarging the generator to
\(H_1(E)\oplus H_1(B)\) includes \(B\) and imposes the new
congruence constraint in (E.2.13b), consistently with E.2.8.
No failure of comparison between places of \(\Q\), or resolution
of the quartic full-category transporter, follows from this
automorphism obstruction.
\end{proof}

\begin{result}{E.2.14}{The common scalar and uniqueness in the prescribed CM fibre}{Proved}
For a full geometric abelian integral comparison as in E.2.9,
the units \(\lambda_A\) are one common
\(\lambda\in\widehat\Z^\times\), over all positive-dimensional
abelian varieties and all Artin-selected embedding blocks.
Fix \(\sigma\in\Gal(\Qb/F_{\mathrm{cm}})\) and prescribe its
particular CM comparison \(\alpha\) from E.2.6. There is at most
one extension of this prescribed comparison to the full geometric
abelian integral diagram. Every such extension satisfies
\[
 T_A\otimes\widehat\Z=\rho_A(\sigma)
 \qquad\hbox{for every positive-dimensional abelian variety }A.
 \tag{E.2.14a}
\]
\end{result}
\begin{proof}
Apply E.2.9 to \(A\times B\) and include the two factor
injections and projections in the diagram. Naturality gives
\[
 T_{A\times B}=T_A\oplus T_B,\qquad
 \rho_{A\times B}(\sigma)=\rho_A(\sigma)\oplus\rho_B(\sigma).
\]
Restricting the scalar equality for \(A\times B\) to each
positive-rank direct summand forces
\(\lambda_A=\lambda_{A\times B}=\lambda_B\).
For varieties over a finite extension of the original field, use
the Artin-module decomposition and selected primitive idempotents
of E.2.3/E.2.9; these preserve precisely the same factor maps.
This proves the common-scalar assertion.

Choose one CM elliptic curve \(C\) in the prescribed subdiagram.
The exact \(s=1\) formula in the main theorem of CM, used in
E.2.6, gives
\[
 \alpha_C\otimes\widehat\Z=\rho_C(\sigma).
\]
The common scalar on this nonzero factor is therefore one.
This proves (E.2.14a) on every factor.

Integral matrices inject into their profinite completions, so
(E.2.14a) determines every \(H_1\) component uniquely. The
cohomology of each abelian variety is the exterior algebra on
\(H^1\). Tensor operations, the prescribed Artin permutation,
and labelled direct-summand projectors then determine all
components in the geometric abelian diagram. This proves
the asserted uniqueness.

In particular, the rational point constructed in E.2.11 must
satisfy the additional equalities (E.2.14a) at every finite place
to yield this integral extension. The Hasse argument constructs
a point in the prescribed CM fibre; it does not establish its
equality to all the actual \(\ell\)-adic Galois comparison
points. For the main place-pair problem, other extensions
\(\sigma\) of the same field embeddings and other compatible
CM comparisons remain to be examined. The further objects of
the full pure-motive diagram remain to be included as well.
\end{proof}
